\section{\label{sec:dicke-grover}Feasible-subspace ansatz and quantum-walk mixers}

\subsection{\label{subsec:qwoa}QWOA on a feasible graph}

In QWOA, the mixer is generated by a quantum walk on a graph of feasible
solutions~\cite{Marsh2019,Marsh2020}. Let
\begin{equation}
    G_{\F}=(\F,E_{\F})
\end{equation}
be a graph whose vertices are feasible bit strings, and $E_{\mathcal{F}}$ is the set of edges. Its adjacency matrix is
\begin{equation}
    \mathcal{A}_{\F}
    =
    \sum_{(x,y)\in E_{\F}}
    \ket{x}\bra{y}.
    \label{eq:feasible-adjacency}
\end{equation}
The QWOA state is
\begin{equation}
    \ket{\psi_p^{\mathrm{QWOA}}}
    =
    \prod_{\ell=1}^{p}
    e^{-i\beta_\ell \mathcal{A}_{\F}}
    e^{-i\gamma_\ell H_C}
    \ket{s_{\F}},
    \label{eq:qwoa-state}
\end{equation}
with
\begin{equation}
    \ket{s_{\F}}
    =
    \frac{1}{\sqrt{|\F|}}
    \sum_{x\in\F}\ket{x}.
    \label{eq:uniform-feasible-state}
\end{equation}
Since $\mathcal{A}_{\F}$ has support only on $\mathcal{H}_{\F}$, the mixer preserves
feasibility.

For fixed cardinality, a natural feasible graph is the Johnson graph
$J(n,k)$. Its vertices are the $k$-element subsets of $\{1,\ldots,n\}$, or
equivalently the bit strings of Hamming weight $k$. Two vertices are adjacent
when the corresponding subsets differ by one element. In the bit-string
representation, this is the same as Hamming distance two. Johnson-graph mixers are natural for
constraints that fix Hamming weight and are closely related to constraint-preserving mixer
constructions in the quantum alternating operator ansatz~\cite{Cook2019}.

Another possible graph is the complete graph $K_{\F}$ on the feasible set. Its
adjacency matrix is
\begin{equation}
    \mathcal{A}_{K_{\F}}
    =
    \sum_{\substack{x,y\in\F\\x\neq y}}
    (\ket{x}\bra{y} + \ket{y}\bra{x}).
    \label{eq:complete-adjacency}
\end{equation}
This mixer connects every feasible portfolio to every other feasible
portfolio.

A natural extension is to keep the Johnson support but assign a different
hopping amplitude to each feasible asset exchange. We define the weighted
Johnson mixer as
\begin{equation}
    \mathcal{A}_{J,w}
    =
    \sum_{(x,y)\in E_J}
    w_{xy}
    \left(
    \ket{x}\bra{y}+\ket{y}\bra{x}
    \right),
    \label{eq:weighted-johnson}
\end{equation}
where $E_J$ is the edge set of the Johnson graph. In the numerical benchmark we
use the cost-similarity weights
\begin{equation}
    w_{xy}
    =
    \exp\left[-\alpha\frac{|C(x)-C(y)|}{\sigma_C}\right],
    \label{eq:cost-weighted-johnson}
\end{equation}
where $\sigma_C$ is the standard deviation of the cost values over the feasible
set \gabriel{Don't we need full access to the feasible set to compute this? If so, this is computationally impossible for reasonably-sized instances, and should only be used in theoretical analyses}. The parameter $\alpha$ controls how strongly transitions between
portfolios with different objective values are suppressed. The case
$\alpha=0$ recovers the unweighted Johnson graph. This extension keeps the
same feasible local moves but makes the quantum walk sensitive to the
cost-landscape geometry.

\subsection{\label{subsec:dicke-grover}Dicke-state ansatz with Grover mixer}

The Dicke state with \(k\) excitations on \(n\) qubits is
\begin{equation}
    \ket{D_k^n}
    =
    \frac{1}{\sqrt{\binom{n}{k}}}
    \sum_{\substack{x\in\{0,1\}^{n}\\ |x|=k}}
    \ket{x}.
    \label{eq:dicke-state}
\end{equation}
For the feasible set \(\mathcal{F}_{n,k}\),
\begin{equation}
    \ket{D_k^n}=\ket{s_{\mathcal{F}}},
\end{equation}
where
\begin{equation}
    \ket{s_{\mathcal{F}}}
    =
    \frac{1}{\sqrt{|\mathcal{F}|}}
    \sum_{x\in\mathcal{F}}\ket{x}.
\end{equation}
Thus Dicke-state preparation gives the QWOA initial state for the
fixed-cardinality problem~\cite{Bartschi2019}.

The feasible Grover mixer is the reflection
\begin{equation}
    G_{\F}
    =
    2\ket{s_{\F}}\bra{s_{\F}}-I_{\F},
    \label{eq:grover-reflection}
\end{equation}
where $I_{\F}$ is the identity on $\mathcal{H}_{\F}$. Grover-type mixers have
been studied as constraint-preserving mixers for QAOA-type algorithms
~\cite{Bartschi2020,Bridi2024}. We now relate this mixer to complete-graph
QWOA.

\begin{proposition}
\label{prop:grover-complete-qwoa}
For the fixed-cardinality feasible set $\F$, the Dicke-state plus Grover-mixer
ansatz is equivalent to QWOA on the complete feasible graph $K_{\F}$, up to a
global phase and a rescaling of the mixing parameter.
\end{proposition}

\begin{proof}
 \gabriel{I haven't checked the proof yet. I'll come back to this.}
 Let
\begin{equation}
    M=|\F|=\binom{n}{k}.
\end{equation}
Define the all-ones operator on $\mathcal{H}_{\F}$ by
\begin{equation}
    J_{\F}=\sum_{x,y\in\F}\ket{x}\bra{y}.
\end{equation}
Since
\begin{equation}
    \ket{s_{\F}}\bra{s_{\F}}=\frac{1}{M}J_{\F},
\end{equation}
and
\begin{equation}
    J_{\F}=\mathcal{A}_{K_{\F}}+I_{\F},
\end{equation}
we obtain
\begin{align}
    G_{\F}
    &=
    2\ket{s_{\F}}\bra{s_{\F}}-I_{\F}
    \nonumber\\
    &=
    \frac{2}{M}\mathcal{A}_{K_{\F}}
    +
    \left(\frac{2}{M}-1\right)I_{\F}.
    \label{eq:grover-complete-relation}
\end{align}
Equivalently,
\begin{equation}
    \mathcal{A}_{K_{\F}}
    =
    \frac{M}{2}G_{\F}
    +
    \left(\frac{M}{2}-1\right)I_{\F}.
    \label{eq:complete-grover-relation}
\end{equation}
Therefore
\begin{align}
    e^{-i\beta \mathcal{A}_{K_{\F}}}
    &=
    \exp\left[
    -i\beta
    \left(
    \frac{M}{2}G_{\F}
    +
    \left(\frac{M}{2}-1\right)I_{\F}
    \right)
    \right]
    \nonumber\\
    &=
    e^{-i\beta\left(\frac{M}{2}-1\right)}
    e^{-i(M\beta/2)G_{\F}}.
    \label{eq:qwoa-grover-equivalence}
\end{align}
The first factor is a global phase. Hence complete-graph QWOA is equivalent to
Grover mixing with the parameter identification
\begin{equation}
    \theta=\frac{M\beta}{2}.
\end{equation}

\end{proof}

This result shows that the Dicke-state plus Grover ansatz realizes a specific
QWOA mixer: the complete feasible graph. It does not imply equivalence to
QWOA on other feasible graphs, such as the Johnson graph. Those choices define
different quantum walks and can have different performance. \gabriel{There are also other known QAOA mixers compatible with Dicke states, for example, the XY-mixer, XY-ring, etc. How do these compare to the QWOA on the Johnson graph?}

\gabriel{In terms of quantum circuits, the Johnson graph QWOA is very vague here. I think it would be nice to describe the circuit construction here. For example, I have no idea what's the depth of the QWOA circuit with each mixer, and this is very relevant if we want to run things on current quantum hardware.}